\chapter{Espacios Vectoriales Booleanos y Códigos de Hamming}

\section{El Espacio Vectorial Binario $\mathbb{B}^n$}

En el Capítulo 6b demostramos que, si bien cualquier álgebra de Boole puede ser interpretada como un anillo, el \textbf{único} álgebra de Boole que tiene la estructura rigurosa de un \textbf{Cuerpo Matemático} (es decir, carente de divisores de cero y donde todo elemento no nulo tiene inverso multiplicativo) es el álgebra bivaluada $\mathbb{B}_2 = \{0, 1\}$. Algebraicamente, este cuerpo es exactamente el cuerpo de Galois $\mathbb{F}_2$.

Dado que la condición insoslayable para construir un espacio vectorial lineal es operar sobre un cuerpo de escalares, deducimos que los únicos espacios vectoriales puramente booleanos que pueden existir deben tener como conjunto base a $\mathbb{F}_2$.

\begin{definicion}{El Espacio Vectorial $\mathbb{B}^n$}{espacio_bn}
Se define el espacio vectorial booleano $\mathbb{B}^n$ como el conjunto de todas las $n$-tuplas (vectores de $n$ bits) cuyos elementos pertenecen a $\mathbb{F}_2$. 
Las dos operaciones que dotan al conjunto de estructura de espacio vectorial son:
\begin{enumerate}
    \item \textbf{Suma vectorial:} Se define como la operación XOR ($\oplus$) aplicada bit a bit entre dos vectores.
    \[ \vec{u} \oplus \vec{v} = (u_1 \oplus v_1, u_2 \oplus v_2, \ldots, u_n \oplus v_n) \]
    \item \textbf{Producto por escalar:} Se define como la operación AND ($\wedge$) entre un escalar booleano $k \in \mathbb{F}_2$ y cada elemento del vector.
    \[ k \wedge \vec{v} = (k \wedge v_1, k \wedge v_2, \ldots, k \wedge v_n) \]
\end{enumerate}
\end{definicion}

Al operar sobre $\mathbb{F}_2$, este espacio vectorial hereda directamente la característica 2 de su cuerpo base. Esto significa que \textbf{todo vector es su propio inverso aditivo}:
\[ \vec{v} \oplus \vec{v} = \vec{0} \]
lo cual simplifica extraordinariamente el cálculo matricial (sumar y restar vectores es exactamente la misma operación).

\section{Métrica: Peso y Distancia de Hamming}

Para que este espacio abstracto tenga utilidad práctica en ingeniería (en concreto, para analizar cómo de parecidos son dos mensajes digitales), necesitamos dotarlo de una métrica que nos permita medir "distancias" entre vectores.

\begin{definicion}{Peso de Hamming}{peso_hamming}
El \textbf{peso de Hamming} de un vector $\vec{v}$, denotado como $w(\vec{v})$, es el número de componentes no nulas (número de unos) que contiene.
\end{definicion}

\begin{definicion}{Distancia de Hamming}{distancia_hamming}
La \textbf{distancia de Hamming} entre dos vectores $\vec{u}$ y $\vec{v}$, denotada como $d(\vec{u}, \vec{v})$, es el número de posiciones en las que difieren. Operacionalmente, coincide con el peso de Hamming de su suma vectorial:
\[ d(\vec{u}, \vec{v}) = w(\vec{u} \oplus \vec{v}) \]
\end{definicion}
La función $d$ cumple todas las propiedades matemáticas de una métrica (no negatividad, identidad de los indiscernibles, simetría y desigualdad triangular).

\section{Códigos de Corrección de Errores (Hamming)}

La aplicación más brillante de dotar a las cadenas de bits de una estructura de espacio vectorial sobre $\mathbb{F}_2$ es la invención de los \textbf{códigos correctores de errores}. 

En un canal de comunicación ruidoso, un vector $\vec{v}$ enviado puede sufrir corrupciones (cambios de 0 a 1 o viceversa), recibiéndose un vector diferente $\vec{r}$. Richard Hamming propuso solucionar esto no usando todo el espacio vectorial $\mathbb{B}^n$, sino limitando los mensajes válidos a un subespacio vectorial más pequeño y controlado.

\subsection{El Subespacio Código y la Matriz de Paridad}
Un código lineal por bloques de longitud $n$ y dimensión $k$ se define matemáticamente como un \textbf{subespacio vectorial} $C \subset \mathbb{B}^n$ de dimensión $k$. Todo subespacio vectorial puede definirse como el núcleo (kernel) de una transformación lineal, representada por una matriz llamada \textbf{Matriz de Paridad ($H$)} de dimensiones $(n-k) \times n$.

\begin{teorema}{Validación del Síndrome}{sindrome_hamming}
Un vector recibido $\vec{r}$ es una palabra código válida (pertenece al subespacio código $C$) si y solo si su producto por la matriz de paridad $H$ (utilizando aritmética en $\mathbb{F}_2$) da como resultado el vector nulo. A este resultado se le denomina \textbf{síndrome} ($\vec{s}$).
\[ \vec{s} = H \cdot \vec{r}^T \]
Si $\vec{s} = \vec{0}$, el vector pertenece al subespacio (no hay errores detectados). Si $\vec{s} \ne \vec{0}$, el vector ha salido del subespacio, lo que indica que se ha corrompido durante la transmisión.
\end{teorema}

Dado que operar matrices sobre $\mathbb{F}_2$ requiere únicamente puertas lógicas XOR y AND, el cálculo del síndrome $\vec{s} = H \cdot \vec{r}^T$ se puede implementar directamente en hardware digital con una eficiencia extrema, constituyendo la base de los modernos sistemas de memoria ECC (Error-Correcting Code).
